\section{Teks Cerpen}
\begin{definition}[Teks Cerpen]
  Suatu bentuk dan hasil pekerjaan seni kreatif yang objek permasalahannya adalah manusia dan kehidupannya, dengan menggunakan bahasa sebagai medianya, serta mengandung amanat atau nilai moral kepada pembaca
\end{definition}
\subsection{Ciri-Ciri Cerpen}
\begin{enumerate}
\item Bentuk tulisan singkat, padat, dan lebih pendek daripada novel.
\item Tulisan kurang dari $10.000$ kata.
\item Sumber cerita dari kehidupan sehari hari.
\item Mengangkat masalah tunggal kehidupan pelaku.
\item Habis dibaca sekali duduk.
\item Tokoh-tokoh nya mengalami konflik sampai pada penyelesaian.
\item Penggunaan kata-katanya khas dan mudah dikenali masyarakat.
\item Meninggalkan kesan mendalam dan efek terhadap ceritanya.
\item Menceritakan suatu kejadian dari terjadinya perkembangan jiwa dan krisis.
\item Beralur tunggal dan lurus.
\item Penokohan singkat, sderhana, dan tidak mendalam.
\end{enumerate}

\subsection{Struktur Cerpen}
\begin{enumerate}
\item \textbf{Abstrak: } Ringkasan atau inti cerita.
\item \textbf{Orientasi: } Pengenalan latar.
\item \textbf{Komplikasi: } Urutan kejadian atau kerumitan yang muncul.
\item \textbf{Evaluasi: } Pemecahan konflik setelah klimaks.
\item \textbf{Resolusi: } Solusi dari berbagai konflik yang dialami tokoh.
\end{enumerate}
